\documentclass[12pt]{article}
\usepackage[utf8]{inputenc}
\usepackage{amsmath}
\usepackage{amssymb}
\usepackage{geometry}
\geometry{a4paper, margin=1in}

\title{Mathematics Test - Solutions}
\author{AI Generated}
\date{\today}

\begin{document}

\maketitle

\section*{Questions}

\subsection*{Question 1}
Solve for $x$: $2x + 5 = 15$

\textbf{Options:}
\begin{itemize}
    \item[A)] $x = 5$
    \item[B)] $x = 10$
    \item[C)] $x = 7.5$
    \item[D)] $x = 20$
\end{itemize}

\subsection*{Question 2}
What is the derivative of $f(x) = x^2 + 3x - 4$?

\textbf{Options:}
\begin{itemize}
    \item[A)] $2x + 3$
    \item[B)] $x + 3$
    \item[C)] $2x$
    \item[D)] $x^2$
\end{itemize}

\subsection*{Question 3}
Evaluate: $\int_0^1 x^2 \, dx$

\textbf{Options:}
\begin{itemize}
    \item[A)] $\frac{1}{2}$
    \item[B)] $\frac{1}{3}$
    \item[C)] $1$
    \item[D)] $\frac{2}{3}$
\end{itemize}

\subsection*{Question 4}
Simplify: $\sqrt{16} + \sqrt{9}$

\textbf{Options:}
\begin{itemize}
    \item[A)] $5$
    \item[B)] $7$
    \item[C)] $25$
    \item[D)] $\sqrt{25}$
\end{itemize}

\subsection*{Question 5}
Solve the quadratic equation: $x^2 - 5x + 6 = 0$

\textbf{Options:}
\begin{itemize}
    \item[A)] $x = 1$ or $x = 6$
    \item[B)] $x = 2$ or $x = 3$
    \item[C)] $x = -2$ or $x = -3$
    \item[D)] $x = 5$ or $x = 1$
\end{itemize}

\newpage

\section*{Solutions}

\subsection*{Solution 1}
\textbf{Answer:} A) $x = 5$

\textbf{Explanation:}
\begin{align*}
    2x + 5 &= 15 \\
    2x &= 15 - 5 \\
    2x &= 10 \\
    x &= 5
\end{align*}

Therefore, $x = 5$ is the correct answer.

\subsection*{Solution 2}
\textbf{Answer:} A) $2x + 3$

\textbf{Explanation:}
Using the power rule and constant rule of differentiation:
\begin{align*}
    f(x) &= x^2 + 3x - 4 \\
    f'(x) &= 2x + 3 + 0 \\
    f'(x) &= 2x + 3
\end{align*}

\subsection*{Solution 3}
\textbf{Answer:} B) $\frac{1}{3}$

\textbf{Explanation:}
Using the power rule for integration:
\begin{align*}
    \int_0^1 x^2 \, dx &= \left[\frac{x^3}{3}\right]_0^1 \\
    &= \frac{1^3}{3} - \frac{0^3}{3} \\
    &= \frac{1}{3}
\end{align*}

\subsection*{Solution 4}
\textbf{Answer:} B) $7$

\textbf{Explanation:}
\begin{align*}
    \sqrt{16} + \sqrt{9} &= 4 + 3 \\
    &= 7
\end{align*}

\subsection*{Solution 5}
\textbf{Answer:} B) $x = 2$ or $x = 3$

\textbf{Explanation:}
Factoring the quadratic equation:
\begin{align*}
    x^2 - 5x + 6 &= 0 \\
    (x - 2)(x - 3) &= 0
\end{align*}

Therefore, $x = 2$ or $x = 3$.

We can verify:
\begin{itemize}
    \item For $x = 2$: $2^2 - 5(2) + 6 = 4 - 10 + 6 = 0$ ✓
    \item For $x = 3$: $3^2 - 5(3) + 6 = 9 - 15 + 6 = 0$ ✓
\end{itemize}

\end{document}
